% AUTO-GENERATED section — do not edit manually
% Place inside the relevant results section, e.g.:
%   % AUTO-GENERATED section — do not edit manually
% Place inside the relevant results section, e.g.:
%   % AUTO-GENERATED section — do not edit manually
% Place inside the relevant results section, e.g.:
%   % AUTO-GENERATED section — do not edit manually
% Place inside the relevant results section, e.g.:
%   \input{sections/healing_procedure}

\subsection{Post-Hoc Feasibility Repair and Healed Objective}
\label{sec:healing}

Several solvers — D-Wave BQM, D-Wave CQM, and the Gurobi decomposed variants —
return solutions that are \emph{infeasible} with respect to one or more
constraints.
To provide a meaningful quality comparison on a common, feasible baseline, we
apply a greedy, \emph{one-shot} repair procedure per solution and report the
resulting \emph{healed objective} alongside (or in place of) the raw objective.
The healed objective is always $\leq$ the raw objective, so it is a
\emph{lower-bound} on the true feasible value the solver would have obtained had
it found a feasible solution.

% -----------------------------------------------------------------------
\subsubsection{Violation Types}
\label{sec:healing:violtypes}

\paragraph{Study~1 (patch benchmark, BQM/CQM).}
The patch-level QUBO/CQM formulation enforces an \emph{at-most-one-crop-per-patch}
constraint.
All observed violations are of this single type: the solver assigns two crops
$c_1, c_2$ to the same patch $p$ (both indicator variables $Y_{p,c_1}$ and
$Y_{p,c_2}$ are set to 1).

\paragraph{Study~2 and Gurobi decomposed (farm benchmark).}
Two constraint types are violated:
\begin{enumerate}
  \item \textbf{One-crop-per-farm} ($\sum_c Y_{f,c} \leq 1$): analogous to the
        patch case — multiple crops are assigned to a single farm.
  \item \textbf{Food-group coverage} ($\sum_{f} Y_{f,c} \geq 1$ for each food
        group): at least one farm must be assigned to a representative crop of a
        specified food group; the solver fails to satisfy one or more of these
        coverage requirements.
\end{enumerate}

% -----------------------------------------------------------------------
\subsubsection{Repair for One-Crop-per-Patch/Farm Violations}
\label{sec:healing:onecrop}

Let $n$ denote the total number of patches (farms), $A_{\text{tot}}$ the total
cultivable area, and $A_p = A_{\text{tot}} / n$ the uniform area per patch
(farm).
Define the normalised per-crop contribution on a violating patch/farm as
\begin{equation}
  \delta_{p,c} \;=\; b_c \cdot \frac{A_p}{A_{\text{tot}}} ,
  \label{eq:patch_contrib}
\end{equation}
where $b_c \geq 0$ is the benefit coefficient of crop $c$ (from the
rotation-benefit matrix).

For each violating patch/farm $p$ with assigned crops $\mathcal{C}_p$
($|\mathcal{C}_p| \geq 2$), we \textbf{remove the minimum-benefit crop}:
\begin{equation}
  \text{gain}_p \;=\; \min_{c\,\in\,\mathcal{C}_p}\; \delta_{p,c} .
  \label{eq:onecrop_gain}
\end{equation}
The healed objective is then
\begin{equation}
  \tilde{f} \;=\; f_{\text{raw}} - \sum_{p \in \mathcal{V}} \text{gain}_p ,
  \label{eq:healed_onecrop}
\end{equation}
where $\mathcal{V}$ is the set of violating patches/farms.
This repair is exact: removing one crop from patch $p$ leaves exactly one
assigned crop, fully restoring feasibility for that patch.

For the \textbf{Gurobi decomposed} variants,
$\tilde{f}$ is pre-computed by the benchmark script and stored as
\texttt{healed\_objective} in
\texttt{solver\_comparison\_results\_hard.json}.  The repair is performed
analytically during the benchmark run using the same procedure above.

For the \textbf{D-Wave BQM/CQM} (Study~1), the repair is applied
post-hoc from the stored \texttt{solution\_plantations} dictionary
$\{Y_{p,c}\}$ by identifying $\mathcal{C}_p = \{c : Y_{p,c} \geq 0.5\}$ for
each violating patch and applying \cref{eq:healed_onecrop}.

% -----------------------------------------------------------------------
\subsubsection{Repair for Food-Group Coverage Violations}
\label{sec:healing:foodgroup}

A food-group constraint of the form
$\sum_{f} Y_{f,g} \geq g_{\min}$
is violated when the number of farms planting a crop from group $g$ falls short
by $\Delta_g = g_{\min} - \sum_f Y_{f,g} > 0$ farms.

The repair swaps $\Delta_g$ existing farm assignments to a representative crop
of group $g$.  We assume the representative crop has benefit $b_g \approx 0$
(starchy-staple crops, which constitute the majority of violated food groups),
so the swap cost equals the full contribution of the farm that is re-assigned.
We choose the \textbf{minimum-benefit farm} to swap, minimising the objective
loss:
\begin{equation}
  \text{gain}_g \;=\; \Delta_g \cdot
    \min_{f\,\in\,\mathcal{F}_{\text{active}}} \;
    \sum_{c\,\in\,\mathcal{C}_f} b_c \cdot \frac{A_f}{A_{\text{tot}}} ,
  \label{eq:foodgroup_gain}
\end{equation}
where $\mathcal{F}_{\text{active}}$ is the set of farms not already assigned to
group $g$.
The total healed objective when both violation types are present is
\begin{equation}
  \tilde{f} \;=\; f_{\text{raw}}
    - \sum_{p\,\in\,\mathcal{V}} \text{gain}_p
    - \sum_{g\,\in\,\mathcal{G}} \text{gain}_g .
  \label{eq:healed_full}
\end{equation}

% -----------------------------------------------------------------------
\subsubsection{Interpretation}
\label{sec:healing:interp}

The healed objective $\tilde{f}$ provides a pessimistic (lower-bound) estimate
of the feasible solution quality.
A large gap $f_{\text{raw}} - \tilde{f}$ signals that the solver's objective
inflation is substantially driven by constraint violations: the reported
$f_{\text{raw}}$ is not a valid crop-allocation benefit.
Conversely, a small gap — as observed for most QPU methods and all Gurobi
decomposed variants beyond $n = 25$ farms — indicates that the solution is
\emph{nearly} feasible and only minor repairs are needed.

In the figures, raw objectives are shown as solid lines and healed objectives
as dotted lines of matching colour.  Scales at which raw and healed coincide
(no violations) are omitted from the dotted series.


\subsection{Post-Hoc Feasibility Repair and Healed Objective}
\label{sec:healing}

Several solvers — D-Wave BQM, D-Wave CQM, and the Gurobi decomposed variants —
return solutions that are \emph{infeasible} with respect to one or more
constraints.
To provide a meaningful quality comparison on a common, feasible baseline, we
apply a greedy, \emph{one-shot} repair procedure per solution and report the
resulting \emph{healed objective} alongside (or in place of) the raw objective.
The healed objective is always $\leq$ the raw objective, so it is a
\emph{lower-bound} on the true feasible value the solver would have obtained had
it found a feasible solution.

% -----------------------------------------------------------------------
\subsubsection{Violation Types}
\label{sec:healing:violtypes}

\paragraph{Study~1 (patch benchmark, BQM/CQM).}
The patch-level QUBO/CQM formulation enforces an \emph{at-most-one-crop-per-patch}
constraint.
All observed violations are of this single type: the solver assigns two crops
$c_1, c_2$ to the same patch $p$ (both indicator variables $Y_{p,c_1}$ and
$Y_{p,c_2}$ are set to 1).

\paragraph{Study~2 and Gurobi decomposed (farm benchmark).}
Two constraint types are violated:
\begin{enumerate}
  \item \textbf{One-crop-per-farm} ($\sum_c Y_{f,c} \leq 1$): analogous to the
        patch case — multiple crops are assigned to a single farm.
  \item \textbf{Food-group coverage} ($\sum_{f} Y_{f,c} \geq 1$ for each food
        group): at least one farm must be assigned to a representative crop of a
        specified food group; the solver fails to satisfy one or more of these
        coverage requirements.
\end{enumerate}

% -----------------------------------------------------------------------
\subsubsection{Repair for One-Crop-per-Patch/Farm Violations}
\label{sec:healing:onecrop}

Let $n$ denote the total number of patches (farms), $A_{\text{tot}}$ the total
cultivable area, and $A_p = A_{\text{tot}} / n$ the uniform area per patch
(farm).
Define the normalised per-crop contribution on a violating patch/farm as
\begin{equation}
  \delta_{p,c} \;=\; b_c \cdot \frac{A_p}{A_{\text{tot}}} ,
  \label{eq:patch_contrib}
\end{equation}
where $b_c \geq 0$ is the benefit coefficient of crop $c$ (from the
rotation-benefit matrix).

For each violating patch/farm $p$ with assigned crops $\mathcal{C}_p$
($|\mathcal{C}_p| \geq 2$), we \textbf{remove the minimum-benefit crop}:
\begin{equation}
  \text{gain}_p \;=\; \min_{c\,\in\,\mathcal{C}_p}\; \delta_{p,c} .
  \label{eq:onecrop_gain}
\end{equation}
The healed objective is then
\begin{equation}
  \tilde{f} \;=\; f_{\text{raw}} - \sum_{p \in \mathcal{V}} \text{gain}_p ,
  \label{eq:healed_onecrop}
\end{equation}
where $\mathcal{V}$ is the set of violating patches/farms.
This repair is exact: removing one crop from patch $p$ leaves exactly one
assigned crop, fully restoring feasibility for that patch.

For the \textbf{Gurobi decomposed} variants,
$\tilde{f}$ is pre-computed by the benchmark script and stored as
\texttt{healed\_objective} in
\texttt{solver\_comparison\_results\_hard.json}.  The repair is performed
analytically during the benchmark run using the same procedure above.

For the \textbf{D-Wave BQM/CQM} (Study~1), the repair is applied
post-hoc from the stored \texttt{solution\_plantations} dictionary
$\{Y_{p,c}\}$ by identifying $\mathcal{C}_p = \{c : Y_{p,c} \geq 0.5\}$ for
each violating patch and applying \cref{eq:healed_onecrop}.

% -----------------------------------------------------------------------
\subsubsection{Repair for Food-Group Coverage Violations}
\label{sec:healing:foodgroup}

A food-group constraint of the form
$\sum_{f} Y_{f,g} \geq g_{\min}$
is violated when the number of farms planting a crop from group $g$ falls short
by $\Delta_g = g_{\min} - \sum_f Y_{f,g} > 0$ farms.

The repair swaps $\Delta_g$ existing farm assignments to a representative crop
of group $g$.  We assume the representative crop has benefit $b_g \approx 0$
(starchy-staple crops, which constitute the majority of violated food groups),
so the swap cost equals the full contribution of the farm that is re-assigned.
We choose the \textbf{minimum-benefit farm} to swap, minimising the objective
loss:
\begin{equation}
  \text{gain}_g \;=\; \Delta_g \cdot
    \min_{f\,\in\,\mathcal{F}_{\text{active}}} \;
    \sum_{c\,\in\,\mathcal{C}_f} b_c \cdot \frac{A_f}{A_{\text{tot}}} ,
  \label{eq:foodgroup_gain}
\end{equation}
where $\mathcal{F}_{\text{active}}$ is the set of farms not already assigned to
group $g$.
The total healed objective when both violation types are present is
\begin{equation}
  \tilde{f} \;=\; f_{\text{raw}}
    - \sum_{p\,\in\,\mathcal{V}} \text{gain}_p
    - \sum_{g\,\in\,\mathcal{G}} \text{gain}_g .
  \label{eq:healed_full}
\end{equation}

% -----------------------------------------------------------------------
\subsubsection{Interpretation}
\label{sec:healing:interp}

The healed objective $\tilde{f}$ provides a pessimistic (lower-bound) estimate
of the feasible solution quality.
A large gap $f_{\text{raw}} - \tilde{f}$ signals that the solver's objective
inflation is substantially driven by constraint violations: the reported
$f_{\text{raw}}$ is not a valid crop-allocation benefit.
Conversely, a small gap — as observed for most QPU methods and all Gurobi
decomposed variants beyond $n = 25$ farms — indicates that the solution is
\emph{nearly} feasible and only minor repairs are needed.

In the figures, raw objectives are shown as solid lines and healed objectives
as dotted lines of matching colour.  Scales at which raw and healed coincide
(no violations) are omitted from the dotted series.


\subsection{Post-Hoc Feasibility Repair and Healed Objective}
\label{sec:healing}

Several solvers — D-Wave BQM, D-Wave CQM, and the Gurobi decomposed variants —
return solutions that are \emph{infeasible} with respect to one or more
constraints.
To provide a meaningful quality comparison on a common, feasible baseline, we
apply a greedy, \emph{one-shot} repair procedure per solution and report the
resulting \emph{healed objective} alongside (or in place of) the raw objective.
The healed objective is always $\leq$ the raw objective, so it is a
\emph{lower-bound} on the true feasible value the solver would have obtained had
it found a feasible solution.

% -----------------------------------------------------------------------
\subsubsection{Violation Types}
\label{sec:healing:violtypes}

\paragraph{Study~1 (patch benchmark, BQM/CQM).}
The patch-level QUBO/CQM formulation enforces an \emph{at-most-one-crop-per-patch}
constraint.
All observed violations are of this single type: the solver assigns two crops
$c_1, c_2$ to the same patch $p$ (both indicator variables $Y_{p,c_1}$ and
$Y_{p,c_2}$ are set to 1).

\paragraph{Study~2 and Gurobi decomposed (farm benchmark).}
Two constraint types are violated:
\begin{enumerate}
  \item \textbf{One-crop-per-farm} ($\sum_c Y_{f,c} \leq 1$): analogous to the
        patch case — multiple crops are assigned to a single farm.
  \item \textbf{Food-group coverage} ($\sum_{f} Y_{f,c} \geq 1$ for each food
        group): at least one farm must be assigned to a representative crop of a
        specified food group; the solver fails to satisfy one or more of these
        coverage requirements.
\end{enumerate}

% -----------------------------------------------------------------------
\subsubsection{Repair for One-Crop-per-Patch/Farm Violations}
\label{sec:healing:onecrop}

Let $n$ denote the total number of patches (farms), $A_{\text{tot}}$ the total
cultivable area, and $A_p = A_{\text{tot}} / n$ the uniform area per patch
(farm).
Define the normalised per-crop contribution on a violating patch/farm as
\begin{equation}
  \delta_{p,c} \;=\; b_c \cdot \frac{A_p}{A_{\text{tot}}} ,
  \label{eq:patch_contrib}
\end{equation}
where $b_c \geq 0$ is the benefit coefficient of crop $c$ (from the
rotation-benefit matrix).

For each violating patch/farm $p$ with assigned crops $\mathcal{C}_p$
($|\mathcal{C}_p| \geq 2$), we \textbf{remove the minimum-benefit crop}:
\begin{equation}
  \text{gain}_p \;=\; \min_{c\,\in\,\mathcal{C}_p}\; \delta_{p,c} .
  \label{eq:onecrop_gain}
\end{equation}
The healed objective is then
\begin{equation}
  \tilde{f} \;=\; f_{\text{raw}} - \sum_{p \in \mathcal{V}} \text{gain}_p ,
  \label{eq:healed_onecrop}
\end{equation}
where $\mathcal{V}$ is the set of violating patches/farms.
This repair is exact: removing one crop from patch $p$ leaves exactly one
assigned crop, fully restoring feasibility for that patch.

For the \textbf{Gurobi decomposed} variants,
$\tilde{f}$ is pre-computed by the benchmark script and stored as
\texttt{healed\_objective} in
\texttt{solver\_comparison\_results\_hard.json}.  The repair is performed
analytically during the benchmark run using the same procedure above.

For the \textbf{D-Wave BQM/CQM} (Study~1), the repair is applied
post-hoc from the stored \texttt{solution\_plantations} dictionary
$\{Y_{p,c}\}$ by identifying $\mathcal{C}_p = \{c : Y_{p,c} \geq 0.5\}$ for
each violating patch and applying \cref{eq:healed_onecrop}.

% -----------------------------------------------------------------------
\subsubsection{Repair for Food-Group Coverage Violations}
\label{sec:healing:foodgroup}

A food-group constraint of the form
$\sum_{f} Y_{f,g} \geq g_{\min}$
is violated when the number of farms planting a crop from group $g$ falls short
by $\Delta_g = g_{\min} - \sum_f Y_{f,g} > 0$ farms.

The repair swaps $\Delta_g$ existing farm assignments to a representative crop
of group $g$.  We assume the representative crop has benefit $b_g \approx 0$
(starchy-staple crops, which constitute the majority of violated food groups),
so the swap cost equals the full contribution of the farm that is re-assigned.
We choose the \textbf{minimum-benefit farm} to swap, minimising the objective
loss:
\begin{equation}
  \text{gain}_g \;=\; \Delta_g \cdot
    \min_{f\,\in\,\mathcal{F}_{\text{active}}} \;
    \sum_{c\,\in\,\mathcal{C}_f} b_c \cdot \frac{A_f}{A_{\text{tot}}} ,
  \label{eq:foodgroup_gain}
\end{equation}
where $\mathcal{F}_{\text{active}}$ is the set of farms not already assigned to
group $g$.
The total healed objective when both violation types are present is
\begin{equation}
  \tilde{f} \;=\; f_{\text{raw}}
    - \sum_{p\,\in\,\mathcal{V}} \text{gain}_p
    - \sum_{g\,\in\,\mathcal{G}} \text{gain}_g .
  \label{eq:healed_full}
\end{equation}

% -----------------------------------------------------------------------
\subsubsection{Interpretation}
\label{sec:healing:interp}

The healed objective $\tilde{f}$ provides a pessimistic (lower-bound) estimate
of the feasible solution quality.
A large gap $f_{\text{raw}} - \tilde{f}$ signals that the solver's objective
inflation is substantially driven by constraint violations: the reported
$f_{\text{raw}}$ is not a valid crop-allocation benefit.
Conversely, a small gap — as observed for most QPU methods and all Gurobi
decomposed variants beyond $n = 25$ farms — indicates that the solution is
\emph{nearly} feasible and only minor repairs are needed.

In the figures, raw objectives are shown as solid lines and healed objectives
as dotted lines of matching colour.  Scales at which raw and healed coincide
(no violations) are omitted from the dotted series.


\subsection{Post-Hoc Feasibility Repair and Healed Objective}
\label{sec:healing}

Several solvers — D-Wave BQM, D-Wave CQM, and the Gurobi decomposed variants —
return solutions that are \emph{infeasible} with respect to one or more
constraints.
To provide a meaningful quality comparison on a common, feasible baseline, we
apply a greedy, \emph{one-shot} repair procedure per solution and report the
resulting \emph{healed objective} alongside (or in place of) the raw objective.
The healed objective is always $\leq$ the raw objective, so it is a
\emph{lower-bound} on the true feasible value the solver would have obtained had
it found a feasible solution.

% -----------------------------------------------------------------------
\subsubsection{Violation Types}
\label{sec:healing:violtypes}

\paragraph{Study~1 (patch benchmark, BQM/CQM).}
The patch-level QUBO/CQM formulation enforces an \emph{at-most-one-crop-per-patch}
constraint.
All observed violations are of this single type: the solver assigns two crops
$c_1, c_2$ to the same patch $p$ (both indicator variables $Y_{p,c_1}$ and
$Y_{p,c_2}$ are set to 1).

\paragraph{Study~2 and Gurobi decomposed (farm benchmark).}
Two constraint types are violated:
\begin{enumerate}
  \item \textbf{One-crop-per-farm} ($\sum_c Y_{f,c} \leq 1$): analogous to the
        patch case — multiple crops are assigned to a single farm.
  \item \textbf{Food-group coverage} ($\sum_{f} Y_{f,c} \geq 1$ for each food
        group): at least one farm must be assigned to a representative crop of a
        specified food group; the solver fails to satisfy one or more of these
        coverage requirements.
\end{enumerate}

% -----------------------------------------------------------------------
\subsubsection{Repair for One-Crop-per-Patch/Farm Violations}
\label{sec:healing:onecrop}

Let $n$ denote the total number of patches (farms), $A_{\text{tot}}$ the total
cultivable area, and $A_p = A_{\text{tot}} / n$ the uniform area per patch
(farm).
Define the normalised per-crop contribution on a violating patch/farm as
\begin{equation}
  \delta_{p,c} \;=\; b_c \cdot \frac{A_p}{A_{\text{tot}}} ,
  \label{eq:patch_contrib}
\end{equation}
where $b_c \geq 0$ is the benefit coefficient of crop $c$ (from the
rotation-benefit matrix).

For each violating patch/farm $p$ with assigned crops $\mathcal{C}_p$
($|\mathcal{C}_p| \geq 2$), we \textbf{remove the minimum-benefit crop}:
\begin{equation}
  \text{gain}_p \;=\; \min_{c\,\in\,\mathcal{C}_p}\; \delta_{p,c} .
  \label{eq:onecrop_gain}
\end{equation}
The healed objective is then
\begin{equation}
  \tilde{f} \;=\; f_{\text{raw}} - \sum_{p \in \mathcal{V}} \text{gain}_p ,
  \label{eq:healed_onecrop}
\end{equation}
where $\mathcal{V}$ is the set of violating patches/farms.
This repair is exact: removing one crop from patch $p$ leaves exactly one
assigned crop, fully restoring feasibility for that patch.

For the \textbf{Gurobi decomposed} variants,
$\tilde{f}$ is pre-computed by the benchmark script and stored as
\texttt{healed\_objective} in
\texttt{solver\_comparison\_results\_hard.json}.  The repair is performed
analytically during the benchmark run using the same procedure above.

For the \textbf{D-Wave BQM/CQM} (Study~1), the repair is applied
post-hoc from the stored \texttt{solution\_plantations} dictionary
$\{Y_{p,c}\}$ by identifying $\mathcal{C}_p = \{c : Y_{p,c} \geq 0.5\}$ for
each violating patch and applying \cref{eq:healed_onecrop}.

% -----------------------------------------------------------------------
\subsubsection{Repair for Food-Group Coverage Violations}
\label{sec:healing:foodgroup}

A food-group constraint of the form
$\sum_{f} Y_{f,g} \geq g_{\min}$
is violated when the number of farms planting a crop from group $g$ falls short
by $\Delta_g = g_{\min} - \sum_f Y_{f,g} > 0$ farms.

The repair swaps $\Delta_g$ existing farm assignments to a representative crop
of group $g$.  We assume the representative crop has benefit $b_g \approx 0$
(starchy-staple crops, which constitute the majority of violated food groups),
so the swap cost equals the full contribution of the farm that is re-assigned.
We choose the \textbf{minimum-benefit farm} to swap, minimising the objective
loss:
\begin{equation}
  \text{gain}_g \;=\; \Delta_g \cdot
    \min_{f\,\in\,\mathcal{F}_{\text{active}}} \;
    \sum_{c\,\in\,\mathcal{C}_f} b_c \cdot \frac{A_f}{A_{\text{tot}}} ,
  \label{eq:foodgroup_gain}
\end{equation}
where $\mathcal{F}_{\text{active}}$ is the set of farms not already assigned to
group $g$.
The total healed objective when both violation types are present is
\begin{equation}
  \tilde{f} \;=\; f_{\text{raw}}
    - \sum_{p\,\in\,\mathcal{V}} \text{gain}_p
    - \sum_{g\,\in\,\mathcal{G}} \text{gain}_g .
  \label{eq:healed_full}
\end{equation}

% -----------------------------------------------------------------------
\subsubsection{Interpretation}
\label{sec:healing:interp}

The healed objective $\tilde{f}$ provides a pessimistic (lower-bound) estimate
of the feasible solution quality.
A large gap $f_{\text{raw}} - \tilde{f}$ signals that the solver's objective
inflation is substantially driven by constraint violations: the reported
$f_{\text{raw}}$ is not a valid crop-allocation benefit.
Conversely, a small gap — as observed for most QPU methods and all Gurobi
decomposed variants beyond $n = 25$ farms — indicates that the solution is
\emph{nearly} feasible and only minor repairs are needed.

In the figures, raw objectives are shown as solid lines and healed objectives
as dotted lines of matching colour.  Scales at which raw and healed coincide
(no violations) are omitted from the dotted series.
